\chapter{2019 年工程流体力学试题（当年科目代码 821）}

\begin{tishi}
\textbf{关于科目代码：}2019 年西华大学该科目的\textbf{考试科目代码为 821}
（原卷封面印"考试科目代码及名称：821 工程流体力学"），与现行的 818 不同。
经渲染原卷目视复核，821 确为原卷所印，\textbf{并非笔误}；
821 与 818 的\textbf{考查内容体系一致}（同为流体力学第 1--8 章加旧大纲/复试范围），
仅代码逐年变动，复习时按同一体系处理即可。

\textbf{试卷说明（原卷原文）：}"试题附在答卷内交回，答案必须写在答题纸上，写在试题纸上无效。"

\textbf{结构：}选择题 10 分（5 小题，每小题 2 分）、判断题 20 分（10 小题，每小题 2 分）、
填空题 30 分（6 小题，每空 3 分，共 10 空）、作图题 7 分（1 小题）、
简答题 16 分（4 小题，6＋4＋3＋3 分）、证明题 17 分（2 小题，7＋10 分）、
计算题 50 分（1、2、3、4 题必做，5、6、7 题\textbf{任选做一道}）。满分 150 分。

\textbf{图示说明：}以下各题〔图示〕均为依据原卷渲染件所作的\textbf{解题用文字描述}，
\textbf{不等于原卷图形}，原卷影印件见本册附录。
\end{tishi}

\section{一、选择题（共 10 分，每小题 2 分）}

\begin{ti}
  \item 气体的粘性主要来自于气体的（\quad）。\hfill\xstarfull{5}\yiju{E4 2026 计算 1（牛顿流体斜板静压强）；E5 题库选择 2（汽油＝牛顿流体）}\nandu{易}
  \begin{choices}
    \item 分子热运动
    \item 分子间的内聚力
    \item 易变形性
    \item 抗拒变形的能力
  \end{choices}

  \item 恒定均匀流的（\quad）。\hfill\xstarfull{3}\yiju{E1 slide16 简答（均匀流的特点）}\nandu{易}
  \begin{choices}
    \item 当地加速度为零
    \item 迁移加速度为零
    \item 向心加速度为零
    \item 合加速度为零
  \end{choices}

  \item 伯努利方程中 $\rho g z+p+\rho\dfrac{\alpha v^{2}}{2}$ 表示（\quad）流体具有的机械能。
  \hfill\xstarfull{5}\yiju{E1 slide16"计算题：伯努利方程"；E1 slide33 与 E2 slide4"第四章＝考试计算题重点"}\nandu{中}
  \begin{choices}
    \item 单位时间
    \item 单位重量
    \item 单位体积
    \item 单位质量
  \end{choices}

  \item 图 1 所示两根完全相同的长管道，只是安装高度不同，两管道的流量关系为（\quad）。
  \hfill\xstarfull{5}\yiju{E5 题库计算 2（水箱垂直管道出流 $Q$ 与管长 $l$）；E1 slide33}\nandu{中}
  \begin{choices}
    \item $Q_1<Q_2$
    \item $Q_1>Q_2$
    \item $Q_1=Q_2$
    \item 不定
  \end{choices}
  \tuyi{原卷图 1：两根完全相同的长管道（等长、等径、等粗糙度）安装高度不同，
  管 1 在上、管 2 在下，两管自左侧引出后各向下弯折；图中 $H$ 为两管弯折段之间的铅垂高差，
  右端画有水面（自由液面符号"$\nabla$"）。原卷未给出管长、管径等数值。
  （原卷影印见本册附录）}

  \item 压强 $p$、密度 $\rho$、长度 $l$、流量 $Q$ 的无量纲组合是（\quad）。
  \hfill\xstarfull{4}\yiju{E1 slide33（量纲转化）；E2 slide5}\nandu{中}
  \begin{choices}
    \item $\dfrac{pl}{\rho Q}$
    \item $\dfrac{pl^{2}}{\rho Q}$
    \item $\dfrac{pl^{4}}{\rho Q^{2}}$
    \item $\dfrac{pl^{2}}{\rho Q^{4}}$
  \end{choices}
\end{ti}

\begin{tishi}
\textbf{第 3 题：}原卷所印表达式为 $\rho g z+p+\rho\dfrac{\alpha v^{2}}{2}$（量纲为压强），
与常见教材"伯努利方程三项之和"的写法略有差异，此处\textbf{照原卷字形转写}，未作改动。

\textbf{第 5 题：}四个选项原卷均为分式（竖排），此处照原卷字形转写；
分母中的密度符号原卷字形与拉丁字母 $p$ 形近（OCR 亦拾作 $p$），
按物理量（密度）统一写作 $\rho$，\textbf{数值与指数均未改动}。
\end{tishi}

\section{二、判断题（正确的打 $\surd$，错误的打 $\times$，每小题 2 分，共计 20 分）}

\begin{ti}
  \item 在圆管流中，层流的断面切应力分布符合抛物线分布规律。\hfill\xstarfull{5}\yiju{E4 2022 单选 10（管径减半压强损失 16 倍）；E4 2015 填空（层流平均速度＝0.5 倍最大速度）}\nandu{中}

  \item 弗劳德数的物理意义是惯性力和重力之比。\hfill\xstarfull{4}\yiju{E4 2014 选择 10；E4 2016 填空 4；E4 2019 判断 2；E4 2022 单选 9；E5 题库选择 12}\nandu{易}

  \item 圆管湍流光滑区的沿程阻力系数 $\lambda$ 只与流动的雷诺数 $\Rey$ 有关。
  \hfill\xstarfull{3}\yiju{E6 教材 5.4；E2 slide5}\nandu{中}

  \item 因为粘性流体的总水头沿程下降，所以水力坡度小于零。\hfill\xstarfull{5}\yiju{E1 slide16 简答"水头损失的分类及定义"；E1 slide33；E2 slide5}\nandu{中}

  \item 按长管计算时可以不考虑局部水头损失和速度水头。\hfill\xstarfull{5}\yiju{E5 题库计算 2（水箱垂直管道出流 $Q$ 与管长 $l$）；E1 slide33}\nandu{易}

  \item 管嘴出流必须满足液体在出口断面满管流出的条件。\hfill\xstarfull{5}\yiju{E4 2022 单选 8（圆柱形直角外管嘴 $H_0\le 9$ m、$l=(3\sim4)d$）；E4 2026 计算 5}\nandu{中}

  \item 在分析有压管路水击现象时，应考虑液体的压缩性。\hfill\xstarfull{4}\yiju{E4 2026 简答 1（水击现象）；E6 教材 5.8}\nandu{中}

  \item 用瑞利法推求物理过程的方程式，待求的量纲指数不能超过 4 个。\hfill\xstarfull{4}\yiju{E4 2016 填空 3；E4 2018 选择 5；E4 2019 选择 5；E4 2020 选择 7}\nandu{中}

  \item 理想均匀流绕流任一静止翼型时，翼型所受的阻力和升力都等于零。\hfill\xstarfull{4}\yiju{E4 2015 计算 8（求驻点与流线）；E2 slide5}\nandu{难}

  \item 平板边界层厚度方向压强不变，即相同 $x$ 坐标边界层内、外的压强相等。
  \hfill\xstarfull{3}\yiju{E6 教材 8.2；E1 slide33}\nandu{中}
\end{ti}

\begin{tishi}
\textbf{第 9 题：}原卷第 9 题与第 10 题之间 OCR 拾得一个孤立字符"Y"；
经原卷渲染件目视核对，该处为扫描噪声，\textbf{第 9、10 题题干均完整无缺损}，
十道判断题的括号均为空（原卷无手写答题痕迹）。
\end{tishi}

\section{三、填空题（共 30 分，每空 3 分）}

\begin{ti}
  \item 欧拉法中运动流体的物理量对时间的导数称为随体导数，由局部导数和位变导数组成。
  任意物理量 $N$ 的随体导数 $\dfrac{\mathrm{d}N}{\mathrm{d}t}=$ \underline{\hspace{3cm}}。
  \hfill\xstarfull{5}\yiju{E4 各年选择/填空（质点导数与加速度）；E1 slide33"公式较多，重点记忆"}\nandu{中}

  \item 已知二维非恒定流场的速度分布为：$u_x=x+t$，$u_y=-y+t$，
  则该流场的流线方程写作 \underline{\hspace{3cm}}。\hfill\xstarfull{5}\yiju{E4 2015 计算 1（流线方程、流动方向判断）}\nandu{中}

  \item 已知平面流动流函数 $\psi=x+y$，则在 $x$、$y$ 轴方向的速度分量
  $u_x=$ \underline{\hspace{2cm}}，$u_y=$ \underline{\hspace{2cm}}，
  在 $x$、$y$ 轴方向的加速度分量 $a_x=$ \underline{\hspace{2cm}}，$a_y=$ \underline{\hspace{2cm}}，
  流动的速度势函数 $\varphi=$ \underline{\hspace{2cm}}。\hfill\xstarfull{5}\yiju{E4 2015 计算（流函数求速度与流线）；E4 2026 计算 4（证明势函数存在性）}\nandu{中}

  \item 水泵的轴功率为 $P$，单位时间内通过水泵的体积流量为 $Q$，水泵的扬程为 $H_m$，
  则该水泵的效率 $\eta=$ \underline{\hspace{3cm}}。\hfill\xstarfull{4}\yiju{E5 题库计算 4（水泵工况点与节流调节）；E6 教材 4.2}\nandu{中}

  \item 如图 2 所示，一高为 $H$、直径为 $D$ 的旋转圆柱容器，充水深度为 $h$，
  如果容器绕 $z$ 轴以等角速度旋转，为使水不溢出，容器的旋转角速度应满足
  \underline{\hspace{3cm}}。\hfill\xstarfull{3}\yiju{E6 教材 2.4；E5 题库}\nandu{难}
  \tuyi{原卷图 2：一圆柱形容器（高 $H$、直径 $D$）绕自身铅垂中心轴 $z$ 以角速度 $\omega$ 旋转；
  容器内充水深度为 $h$（图中 $h$ 自容器底量至初始静止水面，$H$ 自容器底量至容器口缘），
  旋转后水面形成旋转抛物面，图中虚线为初始静止水面，半径标为 $R$（$R=D/2$）。
  容器右上方另标有一个字母（形近 $F$），其含义原卷不清。（原卷影印见本册附录）}

  \item 如图 3 所示，已知两平板间的速度分布为
  $u_x=u_{\max}\left[1-\left(\dfrac{y}{b}\right)^{2}\right]$，式中，$y=0$ 为中心线，
  $y=\pm b$ 为平板所在位置，$u_{\max}$ 为常数，则通过平板的流体单宽流量应为
  \underline{\hspace{3cm}}。\hfill\xstarfull{4}\yiju{E4 2022 单选 4（方管进压缩机求出口密度与质量流量）}\nandu{中}
  \tuyi{原卷图 3：两平行平板间距为 $2b$，以板间中心线为 $x$ 轴（$y=0$），
  两平板分别位于 $y=+b$ 与 $y=-b$；板间速度分布为关于中心线对称的抛物线型，
  中心线处流速最大，标为 $u_{\max}$，流向沿 $x$ 轴正向。（原卷影印见本册附录）}
\end{ti}

\begin{tishi}
\textbf{第 1 题：}原卷该题把"随体导数""局部导数""位变导数"等字样\textbf{直接印在题干中}
（原卷如此，并非填空的答案），空只有 $\dfrac{\mathrm{d}N}{\mathrm{d}t}=$ 一处，
此处照原卷转写、未增删文字。

\textbf{第 2 题：}速度分布原卷 OCR 拾作 $u_x=X+t$，按"二维非恒定流场"题意与
$u_y=-y+t$ 的同构性读作 $u_x=x+t$（\textbf{仅大小写，未改动数值}）。

\textbf{第 3 题：}流函数符号 OCR 拾作 $w$，实为 $\psi$；该题共 5 个空。

\textbf{第 5 题：}原卷图 2 中可见的标注量为 $H$、$h$、$R$、$\omega$（及 $z$ 轴）；
右上角另有一个字形近 $F$ 的字母，其含义与指向\textbf{原卷不清}，请对照原卷影印复核。
\end{tishi}

\section{四、作图题（共 7 分）}

\begin{ti}
  \item 标出图 4 所示盛液容器内 $A$、$B$ 和 $C$ 三点的位置水头、压强水头和测压管水头。
  以图中的 $0$--$0$ 为基准面。\hfill\xstarfull{5}\yiju{E4 2022 单选 3（测压管水头线坡度）、单选 7（虹吸管）；E1 slide33"计算题基础入门"}\nandu{中}
  \tuyi{原卷图 4：一盛液容器，顶部液面（以"$\nabla$"标记）上方的气体压强为 $p_0$；
  容器内标出 $A$（右下部）、$B$（左下部）、$C$（中上部）三点；
  容器右侧壁接一根测压管；基准面 $0$--$0$ 画在容器下方。
  各点与测压管的具体相对位置请对照原卷影印。（原卷影印见本册附录）}
\end{ti}

\section{五、简答题（共 16 分）}

\begin{ti}
  \item 系统和控制体是分析流体运动时经常用到的两个概念，两者的区别是什么？（6 分）
  \hfill\xstarfull{5}\yiju{E4 2015 计算 1（流线方程、流动方向判断）}\nandu{中}

  \item 雷诺实验揭示了沿程损失与流体密切相关，简述不同流态下沿程损失和流速之间的关系。（4 分）
  \hfill\xstarfull{5}\yiju{E4 2022 单选 9（两圆管流态判别）；E4 2026 简答 3（层流与雷诺数关系）}\nandu{易}

  \item 理想均匀流水平向右流动，绕流一个顺时针转动的圆柱体，分析驻点可能出现的位置。（3 分）
  \hfill\xstarfull{4}\yiju{E4 2015 计算 8（求驻点与流线）；E2 slide5}\nandu{难}

  \item 管路的水击可分为直接水击和间接水击，两者哪个的危害更大？为什么？（3 分）
  \hfill\xstarfull{4}\yiju{E4 2026 简答 1（水击现象）；E6 教材 5.8}\nandu{中}
\end{ti}

\section{六、证明题（共 17 分）}

\begin{ti}
  \item 如图 5 所示，蓄水池侧壁装有一直径为 $D$ 的圆形闸门，闸门平面与水面夹角为 $\theta$，
  闸门形心 $C$ 处水深 $h_C$，闸门可绕通过形心 $C$ 的水平轴旋转，
  证明作用于闸门水压力对轴的力矩与形心水深 $h_C$ 无关。（7 分）
  \hfill\xstarfull{5}\yiju{E4 2015 计算（圆柱闸门）；E1 slide33}\nandu{难}
  \tuyi{原卷图 5：蓄水池侧壁上斜装一圆形闸门，闸门平面与水面夹角为 $\theta$，
  图中闸门直径标为 $125\,\mathrm{cm}$；$C$ 为闸门形心（其水深标为 $h_C$），
  $D$ 为形心下方的一点（其水深标为 $h_D$）；$O$--$O$ 为过形心 $C$ 的水平轴，
  $P$ 表示水压力的作用方向。原卷图中 $D$ 点与 $h_D$ 的标注位置请对照影印件。
  （原卷影印见本册附录）}

  \item 如图 6，流量为 $Q$、平均流速为 $v$ 的射流，冲击直立平板后分成两股，
  一股沿板面直泻而下，流量为 $Q_1$，另一股以倾角 $\alpha$ 射出，流量为 $Q_2$，
  两股分流的平均速度均等于 $v$。若不计摩擦力和重力影响，试推证：固定平板所需的外力
  \[
    F=\rho Qv\left(1-\sqrt{\frac{1-Q_1/Q_2}{1+Q_1/Q_2}}\right)
  \]
  （10 分）\hfill\xstarfull{5}\yiju{E4 2015 计算（水平弯管求 $F_x,F_y$）；E4 2026 计算 3（教材习题 4.12）；E1 slide16}\nandu{难}
  \tuyi{原卷图 6：水平射流（流量 $Q$、平均流速 $v$）自左侧冲击一块直立平板，
  来流断面为 $0$--$0$；射流在板面处分成两股，一股沿板面直泻而下（断面 $1$--$1$、流量 $Q_1$、流速 $v$），
  另一股与板面成倾角 $\alpha$ 斜向射出（断面 $2$--$2$、流量 $Q_2$、流速 $v$）。
  （原卷影印见本册附录）}
\end{ti}

\begin{tishi}
\textbf{第 1 题：}原卷图 5 中除 $D$、$\theta$、$h_C$ 外，还标有闸门直径 $125\,\mathrm{cm}$、
形心下方点的水深 $h_D$ 与过形心的水平轴 $O$--$O$；闸门直径数值\textbf{照原卷图字转写}，
请对照原卷影印核实。

\textbf{第 2 题：}原卷所印待证公式经渲染目视核对为
$F=\rho Qv\left(1-\sqrt{\dfrac{1-Q_1/Q_2}{1+Q_1/Q_2}}\right)$（根号覆盖整个分式），
此处照原卷字形转写。原卷 OCR 对该式拾取严重残缺，请对照影印件复核。
\end{tishi}

\section{七、计算题（共 50 分）（1、2、3、4 题必做，5、6、7 题任选做一道）}

\begin{ti}
  \item 如图 7 所示，利用文丘里流量计测量竖直水管中的流量。已知 $d_1=300\,\mathrm{mm}$，
  $d_2=150\,\mathrm{mm}$，水银压差计读数 $\Delta h=20\,\mathrm{mm}$。试确定水流量 $Q$。（10 分）
  \hfill\xstarfull{5}\yiju{E4 2022 单选 6（计算点选取）；E4 2026 计算 2（文丘里流量证明）}\nandu{难}
  \tuyi{原卷图 7：竖直水管上装文丘里流量计，管中收缩喉部断面为 $2$--$2$、
  较宽的断面为 $1$--$1$；两个断面各引出一根细管接到一只 U 形水银压差计上，
  读数 $\Delta h$ 为两管水银面的高差；图中两处直径标注的字形与 $d_1$ 形近，
  按题干数值应对应 $d_1=300\,\mathrm{mm}$（宽断面）与 $d_2=150\,\mathrm{mm}$（喉部），
  图注下标请对照原卷影印复核。（原卷影印见本册附录）}

  \item 如图 8 所示，利用毛细管测定油液粘度，已知毛细管直径 $d=4.0\,\mathrm{mm}$，
  长度 $L=0.5\,\mathrm{m}$，流量 $Q=1.0\,\mathrm{cm^3/s}$ 时，测压管落差 $h=15\,\mathrm{cm}$。
  管中作层流动，求油液的运动粘度。（10 分）\hfill\xstarfull{5}\yiju{E4 2022 单选 10（管径减半压强损失 16 倍）；E4 2015 填空（层流平均速度＝0.5 倍最大速度）}\nandu{中}
  \tuyi{原卷图 8：水平放置的毛细管（直径 $d$、长 $L$）两端接两根竖直的测压管，
  两测压管中液面的高度差为 $h$；油液自毛细管中流出，出口流量为 $Q$。
  （原卷影印见本册附录）}

  \item 已知直径为 $15\,\mathrm{cm}$ 的输油管，流量为 $0.18\,\mathrm{m^3/s}$。
  现用流量为 $0.018\,\mathrm{m^3/s}$ 的水作模型实验，测得 $5\,\mathrm{m}$ 长输水管两端的
  压强水头差 $\dfrac{\Delta p_m}{\rho_m g_m}=5\,\mathrm{cm}$。若模型的管径与原型相同，
  试求 $100\,\mathrm{m}$ 长的输油管两端的压强差 $\dfrac{\Delta p_p}{\rho_p g_p}$？
  （用油柱高表示）（10 分）\hfill\xstarfull{3}\yiju{E4 2014 计算 8（溢流坝模型相似）；E5 题库选择 15}\nandu{难}

  \item 一输水管直径 $d=250\,\mathrm{mm}$，管长 $L=200\,\mathrm{m}$，管壁的切应力
  $\tau=46\,\mathrm{N/m^2}$，求在 $100\,\mathrm{m}$ 长管上的水头损失及在圆管中心和
  $r=100\,\mathrm{mm}$ 处的切应力。（10 分）\hfill\xstarfull{5}\yiju{E4 2022 单选 3；E5 题库选择 10（圆管与方管沿程损失比 0.886）}\nandu{中}

  \item 如图 9，将速度为 $v_\infty$ 的平行于 $x$ 轴的均匀流和在原点强度为 $q$ 的点源叠加，
  求叠加后流场中驻点位置，及经过驻点的流线方程。（10 分）\hfill\xstarfull{4}\yiju{E4 2015 计算 8（求驻点与流线）；E2 slide5}\nandu{中}
  \tuyi{原卷图 9：平面直角坐标系 $Oxy$ 中，沿 $x$ 轴正方向（自左向右）的均匀来流 $v_\infty$，
  在原点 $O$ 处置一强度为 $q$ 的点源；图中画出了叠加后的流线（含经过驻点的分界流线）
  与夹角 $\theta$。（原卷影印见本册附录）}

  \item 滞止参数为 $p_0=10.35\times10^{5}\,\mathrm{Pa}$，$T_0=350\,\mathrm{K}$ 的空气气流
  进入收缩喷管，出口断面的直径 $d=15\,\mathrm{mm}$，若出口外部的环境压强
  $p_b=5\times10^{5}\,\mathrm{Pa}$，试求喷管的质量流量 $Q_m$。
  （空气的气体常数 $R=287\,\mathrm{J/(kg\cdot K)}$，等熵指数 $k=1.4$）（10 分）
  \hfill\xstarmark{X2}\nandu{难}

  \item 使小钢球在油中自由沉降以测定油的粘度。已知油的密度 $\rho=900\,\mathrm{kg/m^3}$，
  小钢球的直径 $d=3\,\mathrm{mm}$，密度 $\rho'=7788\,\mathrm{kg/m^3}$，
  若测得钢球最终的沉降速度 $v=12\,\mathrm{cm/s}$，求油的动力粘度。（10 分）
  \hfill\xstarfull{5}\yiju{E4 2026 计算 1（牛顿流体斜板静压强）；E5 题库选择 2（汽油＝牛顿流体）}\nandu{难}
\end{ti}

\begin{tishi}
\textbf{第 1 题：}原卷图 7 中两处直径标注的字形与 $d_1$ 形近（渲染目视核对亦如此），
按题干给出的 $d_1=300\,\mathrm{mm}$、$d_2=150\,\mathrm{mm}$，
喉部应为 $d_2$、宽断面应为 $d_1$；\textbf{图注下标请对照原卷影印复核}。

\textbf{第 2 题：}原卷流量 OCR 拾作"$1.0\,\mathrm{cm^2/s}$"，
经渲染目视核对，原卷所印为 $Q=1.0\,\mathrm{cm^3/s}$（指数为 3），已按原卷图字转写。

\textbf{第 3 题：}该题公式在原卷 OCR 中严重残缺（仅存"\dots Pm8m\dots Ppgp"等碎片），
经渲染目视核对，原卷所印为"测得 $5\,\mathrm{m}$ 长输水管两端的压强水头差
$\dfrac{\Delta p_m}{\rho_m g_m}=5\,\mathrm{cm}$"，并求 $\dfrac{\Delta p_p}{\rho_p g_p}$
（用油柱高表示）；此处照原卷图字转写，\textbf{数值未改动}。

\textbf{第 4 题：}$100\,\mathrm{m}$ 长管与原卷给出的管长 $L=200\,\mathrm{m}$ 并存，
照原卷字面转写（管长与所求段长的关系请对照原卷影印核实）。

\textbf{第 6 题：}原卷环境压强 OCR 拾作 $p$，经渲染目视核对为 $p_b=5\times10^{5}\,\mathrm{Pa}$
（下标 $b$ 表示环境压强）；滞止压强指数位字形较模糊，照原卷读法转写为 $10.35\times10^{5}\,\mathrm{Pa}$，
\textbf{未改动数值}。本题属\textbf{气体动力学（一元可压缩流动）}内容，
为旧大纲 / 复试范围，现行 818 大纲不要求（考点码 X2）。

\textbf{第 7 题：}原卷钢球密度 OCR 拾作 $\rho'=7788\,\mathrm{kg/m^2}$，
密度单位按物理量改正为 $\mathrm{kg/m^3}$，\textbf{数值 $7788$ 未改动}；
油液密度 $\rho=900\,\mathrm{kg/m^3}$ 与沉降速度 $v=12\,\mathrm{cm/s}$ 照原卷转写。
\end{tishi}
